\iffalse
SimpleTrainSim - application immitating train HUD.
Copyright (C) 2026 Grzegorz Śnieżko

This program is free software: you can redistribute it and/or modify
it under the terms of the GNU General Public License as published by
the Free Software Foundation, either version 3 of the License, or
(at your option) any later version.

This program is distributed in the hope that it will be useful,
but WITHOUT ANY WARRANTY; without even the implied warranty of
MERCHANTABILITY or FITNESS FOR A PARTICULAR PURPOSE.  See the
GNU General Public License for more details.

You should have received a copy of the GNU General Public License
along with this program.  If not, see <https://www.gnu.org/licenses/>.
\fi

\documentclass[12pt,a4paper]{report}

\usepackage[utf8]{inputenc}
\usepackage[T1]{fontenc}
\usepackage[english]{babel}

\usepackage{enumitem} % smaller distance between list items, automatic FR-xxx-000 enumeration
\usepackage{fmtcount} % Provides \padzeroes
\usepackage{geometry}
\usepackage{hyperref}
\hypersetup{
    colorlinks=true,
    linkcolor=black,
    citecolor=black,
    urlcolor=blue
}
\usepackage{setspace}
\usepackage{parskip} % disable paragraph insets
\usepackage{pdflscape}
\usepackage{titlesec} % for altering title headers
\usepackage[titletoc]{appendix} % for word "appendix" in the table of contents
\usepackage{graphicx} % for UML diagrams
\usepackage{verbatim} % for appending GNU GPL3 license
\usepackage{listings} % for programming code
\usepackage{xcolor}

\titleformat{\chapter}[hang]{\normalfont\huge\bfseries}{\thechapter. }{0pt}{\huge}
\titlespacing*{\chapter}{0pt}{-30pt}{20pt}

% C++ code styling from https://www.overleaf.com/learn/latex/Code_listing
\definecolor{codegreen}{rgb}{0,0.6,0}
\definecolor{codegray}{rgb}{0.5,0.5,0.5}
\definecolor{codepurple}{rgb}{0.58,0,0.82}
\definecolor{backcolour}{rgb}{0.95,0.95,0.92}

\lstdefinestyle{mystyle}{
    backgroundcolor=\color{backcolour},   
    commentstyle=\color{codegreen},
    keywordstyle=\color{magenta},
    numberstyle=\tiny\color{codegray},
    stringstyle=\color{codepurple},
    basicstyle=\ttfamily\footnotesize,
    breakatwhitespace=false,         
    breaklines=true,                 
    captionpos=b,                    
    keepspaces=true,                 
    numbers=left,                    
    numbersep=5pt,                  
    showspaces=false,                
    showstringspaces=false,
    showtabs=false,                  
    tabsize=2
}

\lstset{style=mystyle}

% Custom title page
\newcommand{\makecustomtitle}[5]{%
    \newgeometry{margin=1.5in}
    \thispagestyle{empty}
    
    \begin{center}
        \vspace*{\fill}
        
        {\Huge \textbf{#1}} \\[1em]       % Title (28pt)
        {\Large \textbf{#2}} \\[2em]      % Short description (24pt)
        
        {\large Version #3} \\[1.5em]     % Version (16pt)
        
        {\large \today} \\[3em]           % Date
        
        {\large #4}                       % Author (14pt)
        
        \vspace*{\fill}

        {\large #5}                       % License
    \end{center}
    
    \clearpage
    \restoregeometry
    
    \tableofcontents
    \clearpage
}

% Functional and non-functional requirements list auto-numerating IDs
\DeclareRobustCommand{\reqLabel}[2]
{\hspace{16pt}\textbullet\hspace{7pt}\bfseries#1-\padzeroes[3]{\decimal{#2}}:}

\begin{document}

\makecustomtitle
  {SimpleTrainSim} % Title
  {Train HUD}        % Short description
  {0.1}              % Version
  {Grzegorz Śnieżko} % Author
  {GNU GPL v3.0}         % License

\chapter{Domain description}
\section{Product vision}
Simple Qt application designed to simulate realistic train physics and display corresponding data in GUI. GUI is styled after train Head-Up Display (HUD).

\section{Detailed requirements}
    \subsection{Environmental resistance} \label{sec:davis}
    In order to move, the locomotive must overcome environmental drag. 
    Total running resistance ($R$) is calculated using \textbf{Davis equation}:
        \begin{equation}
            R = A + Bv + Cv^2
        \end{equation}
    where 
    \begin{itemize}[noitemsep]
        \item $v$ is the train speed,
        \item $A$ is mechanical rolling resistance,
        \item $B$ is dynamic mechanical forces (e.g., axle friction),
        \item $C$ is aerodynamic drag coefficiency.
    \end{itemize}

    \subsection{Coasting}
    Cutting off the engine power at high speed must let the train roll and gradually slow down because of resistance forces.
    
    \subsection{Bearing friction}
    Starting resistance is modeled higher than dynamic resistance
    due to bearing/axle friction mechanics.

    \subsection{Adhesion and weather}
    Excessive engine force can cause the wheel slip. 
    Bad weather conditions can decrease the amount of force needed for wheels to start slipping.
    
    \subsection{Gravitational influence}
    Gravity influences the train based on the current track gradient.
    It can accelerate or deaccelerate the train.
    
    \subsection{Aerodynamics}
    Wind resistance increases non-linearly with speed. 
    Every train profile holds a value of its aerodynamic efficiency.
    
    \subsection{Safety systems}
    User must periodicaly interact with deadman's switch. 
    Time between safety warnings depends on train model.

\chapter{Functional requirements}
\section{Graphical User Interface (GUI)}
\begin{enumerate}[label=\reqLabel{FR-GUI}*, align=left]
    \item System must render GUI styled after train Head-Up Display (HUD).
    \item The HUD must display following data updated in real time:
    \begin{itemize}
        \item speed in $\mathrm{km/h}$, 
        \item engine force in $\mathrm{kN}$, 
        \item throttle setting in $\%$, 
        \item wheel slip warning, 
        \item braking force in $\mathrm{bars}$, 
        \item brake setting in $\%$, 
        \item track gradient in $\%$, 
        \item acceleration in $\mathrm{m/s^2}$
    \end{itemize} 
    \item Application must start in start menu with following options:
    \begin{itemize}
        \item ride
        \item add route
        \item add train
        \item exit application
    \end{itemize}
    \item \textit{Add route} panel must let the user upload a JSON file with route height points.
    \item \textit{Add train} panel must display train creation wizard. Some fields must allow the user to choose predefined values:
    \begin{itemize}
        \item engine type (electric, combustion with electric drive, combustion with direct drive)
        \item HUD layout
        \item speed indicator design
        \item throttle indicator design
        \item brake indicator design
        \item throttle lever
        \item brake lever
    \end{itemize}
    Some fields will let the user to input their own values:
    \begin{itemize}
        \item locomotive weight
        \item drag coefficiency
        \item engine power
    \end{itemize}
\end{enumerate}

\section{Physics Simulation}
\begin{enumerate}[label=\reqLabel{FR-SIM}*, align=left]
    \item The system must calculate total running resistance at every simulation frame using the Davis equation defined in section~\ref{sec:davis}.
    \item The system must apply bearing friction to resistance calculation as long as the train speed is 0~$\mathrm{km/h}$
    \item The system must simulate coasting behaviour when the train is moving and current engine force is 0. Current coasting speed is altered by current resistance forces and gravitaional forces.
    \item The system must calculate the effects of gravity on train in real time based on current track gradient. Gravity will move the stopped train on a slope if no brakes are applied. During coasting uphill, the gravity will be slowing the train down.
\end{enumerate}

\section{Environmental factors}
\begin{enumerate}[label=\reqLabel{FR-ENV}*, align=left]
    \item Route is made of points that are 100 metres away from each other. Every point must hold the altitude value in metres above mean sea level.
    \item Current weather must affect maximum track adhesion level. Damp track must have lower adhesion. Icy track must have even lower adhesion level. Track must have the leaf coverage percentage which will also affect the adhesion.
    \item Train route must have at least one wind zone. Wind zone will hold two values:
    \begin{itemize}
        \item wind speed in $\mathrm{m/s}$,
        \item wind direction in degrees (0 to 360)
    \end{itemize}
    Route will always point to the north (0 deg). Wind affects total running resistance.
\end{enumerate}

\section{Safety systems}
\begin{enumerate}[label=\reqLabel{FR-SAF}*, align=left]
    \item The train must maintain an active timer for deadman's switch.
    \item The user can reset deadman's switch timer by pressing an appropriate button or keyboard shortcut.
    \item The train must apply an emergency brake and cut the engine power if user doesn't manage to respond in appropriate time.
    \item In order to release an emergency brake, user must:
    \begin{enumerate}
        \item respond to deadman's switch,
        \item set the throttle position to neutral,
        \item set the brake lever to emergency position
    \end{enumerate}
    When the combination is complete user can set the brake to release or neutral position.
\end{enumerate}
    


\chapter{Non-functional requirements}
\section{Performance}
\begin{enumerate}[label=\reqLabel{NFR-PER}*, align=left]
    \item Control input latency should be less than 50ms.
    \item Application must maintain a stable memory footprint without memory leaks.
\end{enumerate}

\section{Reliability}
\begin{enumerate}[label=\reqLabel{NFR-REL}*, align=left]
    \item User must be able to complete the route without crashes and then change it.
    \item Application must exit only on user command.
\end{enumerate}

\section{Usability}
\begin{enumerate}[label=\reqLabel{NFR-USA}*, align=left]
    \item Application must provide GUI with mouse control and keyboard shortcuts.
    \item GUI must meet the WCAG standards at least at the AA level, preferably AAA.
\end{enumerate}

\section{Portability}
\begin{enumerate}[label=\reqLabel{NFR-POR}*, align=left]
    \item Application should run on most popular PC operating systems: Windows, macOS, Linux
\end{enumerate}

\section{Maintainability \& Extensibility}
\begin{enumerate}[label=\reqLabel{NFR-MNT}*, align=left]
    \item Code must be self-documenting with intuitive naming conventions.
    \item Code must follow object-oriented programming rules to ensure modularity.
    \item User must have the ability to add different routes and trains by providing only their data. Data must be persistent.
\end{enumerate}

\chapter{Use Case Diagram}
\section{Actors}
\begin{itemize}
    \item \textbf{Driver} -- controls train through interface
    \item \textbf{Deadman's switch} -- safety system, controls train directly
    \item \textbf{Train} -- physical actor, moves
    \item \textbf{Gravity} -- physical force, affects train speed
    \item \textbf{Environment} -- sum of physical resistance forces
\end{itemize}
\begin{figure}[htbp]
    \centering
    \vspace*{-2cm}
    \hspace*{-2.5cm}
    \includegraphics[width=20cm]{uml/UCD_Ride.png}
    \caption{Use Case Diagram - Riding}
\end{figure}

\newgeometry{margin=-1in}
\begin{landscape}
\titlespacing*{\chapter}{90pt}{60pt}{0pt}
\chapter{Conceptual Class Diagram}
\begin{figure}[htbp]
    \centering
    \includegraphics[width=28.5cm]{uml/conceptualClassDiagram.png}
    \caption{SimpleTrainSim Conceptual Class Diagram}
    \label{fig:conceptual}
\end{figure}


\chapter{Design Class Diagram}
\begin{figure}[htbp]
    \centering
    \includegraphics[width=28.5cm]{uml/DesignCLD_Overview.png}
    \caption{Design CLD - Overview}
\end{figure}

\begin{figure}[htbp]
    \vspace*{-3cm}
    \centering
    \includegraphics[width=28.5cm]{uml/DesignCLD_Train.png}
    \caption{Design CLD - Train}
    \vspace*{-3cm}
\end{figure}

\begin{figure}[htbp]
    \centering
    \includegraphics[width=28.5cm]{uml/DesignCLD_UI.png}
    \caption{Design CLD - UI}
\end{figure}

\begin{figure}[htbp]
    \centering
    \includegraphics[width=28.5cm]{uml/DesignCLD_Route.png}
    \caption{Design CLD - Route}
\end{figure}

\begin{figure}[htbp]
    \centering
    \includegraphics[width=28.5cm]{uml/DesignCLD_Physics.png}
    \caption{Design CLD - Physics}
\end{figure}

\end{landscape}
\restoregeometry
\titlespacing*{\chapter}{0pt}{-30pt}{20pt}

\chapter{Dynamic analysis}
\section{State diagrams}
\begin{figure}[htbp]
    \centering
    \includegraphics[height=15cm]{uml/stateBrake.png}
    \caption{State diagram - brake}
\end{figure}

\begin{figure}[htbp]
    \centering
    \hspace*{-1.5cm}
    \includegraphics[width=18cm]{uml/stateDeadmansSwitch.png}
    \caption{State diagram - deadman's switch}
\end{figure}

\begin{figure}[htbp]
    \centering
    \includegraphics{uml/stateEngine.png}
    \caption{State diagram - engine}
\end{figure}

\section{Sequence diagrams}
\begin{figure}[htbp]
    \centering
    \includegraphics[width=15cm]{uml/sequenceEmergencyBrakeReset.png}
    \caption{Sequence diagram - emergency brake}
    \vspace*{-3cm}
\end{figure}

\vspace*{3cm}
\section{Maths}
\subsection{Engine}
Engine transform power to tractive force based on this formula:
\[F = \frac{P}{|v|}\]
where:
\begin{itemize}
    \item F - tractive force (N)
    \item P - generated engine power (W)
    \item v - absolute train speed (m/s)
\end{itemize}

Because train is given maximum tractive force $F_{max}$, engine power $P$ is set properly to not create a force bigger than maximum tractive force.
\[P = F_{max} * v\]
With formula above train would not generate power at $v=0$. Because of this, train engine assumes the speed doesn't fall below certain value.
\lstinputlisting[language=C++]{code_snippets/engine_min_speed.cpp}

Train has artificial speed limiter that smoothly decreases power output after reaching certain speed.
\[DropFactor = max(0, 1 - \frac{v - v_{StartFade}}{FadeRange})\]
\[P = F_{max} * v * DropFactor\]
where
\begin{itemize}
    \item $v_{StartFade}$ - speed at which engine starts to decrease power
    \item $FadeRange$ - the difference between maximum speed $v_{max}$ and $v_{StartFade}$
\end{itemize}
That means:
\[v_{max} = v_{StartFade} + FadeRange\]

Engine smoothly increases power
\[P = \Delta P_{rate} * \Delta T\]
where $\Delta T$ is the time between function calls taken from function argument:
\begin{lstlisting}
void Engine::calculate_force(..., double deltaTimeInSeconds)
\end{lstlisting}
$\Delta P_{rate}$ is engine variable that tells by how much each function call can increase the power.

Finally, the tractive force that is available to read by other \verb|Q_OBJECTS| is calculated with previously mentioned formula $F = \frac{P}{|v|}$

\subsection{Train}
Train calculates net force as the difference between tractive force (calculated by the engine) and total resistance force.
Total resitance force is calculated from Davis Formula.
Constant values are taken from work \href{https://rosap.ntl.bts.gov/view/dot/79072}{"The Tractive Resistance of Electric Locomotives and Cars"} by W. J. Davis.
\[R = A + B*v + C*v^2\]
\[A = 1.3\frac{lbf}{ton} + \frac{29\frac{lbf}{ton}}{w}\]
$w$ is weight per axle in US tons
\[B = 0.045\frac{lbf}{ton}\]
\[C = \frac{0.0005\frac{lbf}{ton} * a}{w * n}\]
\begin{itemize}
    \item $a$ - frontal area in $sqft$)
    \item $w$ - weight per axle in US tons
    \item $n$ - number of axles
\end{itemize}

Because program data is metric, $\frac{lbf}{ton}$ had to be converted to $\frac{N}{kg}$.
\[1 lbf = 4.44822162 N\]
\[1 (US)Ton = 0.9071847 (Metric)Ton\]
\[1\frac{lbf}{ton} = 4.90332522\frac{N}{t}\]

That gives:
\[A \approx 6.374 + \frac{142.1964}{w * 0.9071847} \approx 6.374 + \frac{128.99}{w}\]
\[B \approx 0.2206\]
After convertin $mph$ to $km/h$ we can use:
\[B \approx 0.1371\]
$C$ needs to be recalculated

Train acceleration is calculated using
\[a = \frac{f}{m}\]
\begin{itemize}
    \item $f$ is net force in N
    \item $m$ is train mass in kg
\end{itemize}

Speed is calculated using
\[\Delta v = a * \Delta t\]
$\Delta t$ is time since last function call

\begin{appendices}

\chapter{GNU General Public License v3.0}
\label{app:license}

This software is released under the \href{https://www.gnu.org/licenses/gpl-3.0.en.html#license-text}{GNU General Public License v3.0}. 
Source code is available on \href{https://github.com/grzeg2010/SimpleTrainSim}{GitHub}.
The full text of the license is included below.

\begin{singlespace}
    \small
    \verbatiminput{../LICENSE}
\end{singlespace}

\end{appendices}
\end{document}